% !TEX program = lualatex
\documentclass[12pt,a4paper]{ltjsarticle}
\usepackage{amsmath,amssymb,amsthm}
\usepackage{geometry}
\geometry{margin=2.5cm}
\usepackage{hyperref}
\hypersetup{colorlinks=true,linkcolor=blue,citecolor=blue,urlcolor=blue}
\usepackage{booktabs}
\usepackage{luatexja-fontspec}
\usepackage{graphicx}

\title{ディレイ回路モデルに内在する対称性の整理}
\author{木原 範昭（Noriaki Kihara）\\
WF System Co., Ltd.\\
大阪大学 基礎工学部 卒業}
\date{2026年4月\quad 研究ノート}

\begin{document}

\maketitle

\begin{abstract}
\end{abstract}

\section{はじめに}

\section{準備：前稿の定義の要約}

\section{対称性の整理}

\subsection{$I$ と $\overline{I}$ の対称性}

前稿~\cite{ref1} \S 2.2 性質5 において，反転 $\overline{I}$ は同一クラス内の操作と定義された．すなわち，$I$ と $\overline{I}$ は同一クラスに属する．

ディレイ回路 $F_D$ の演算規則は，入力情報と内部情報の等価判定（$\mathrm{IN} \neq I$ か否か）のみに依存する．$I$ を $\overline{I}$ で置き換えても，等価判定と演算規則の構造は変わらない．

\textbf{定義（反転対称性）}：ディレイ回路 $F_D$ の演算規則は，$I \to \overline{I}$ の変換に対して不変である．これを反転対称性と呼ぶ．

この性質は，反転が同一クラス内の操作であり $\overline{\overline{I}} = I$ を満たすという定義（\cite{ref1} \S 2.2）から直接導かれる自明な帰結である．

\subsection{スーパークラスの定義}

前稿~\cite{ref1}では，同一クラス内の演算（ディレイ回路 $F_D$ および NOT演算回路）のみを扱った．本節では，異なるクラスに属する情報を統合的に扱うための概念を定義する．

異なるクラスに属する情報の組 $(A, B, C_1, C_2, \ldots, C_m)$ を\textbf{スーパークラス} $\mathcal{I}$ と定義する．子要素の個数は任意（0個を含む）であり，各子要素のクラスと値も制約されない．子要素が0個の場合は $\mathrm{null}$ に相当する．

スーパークラスは以下の性質を満たすものとする．

\paragraph{性質1：順序非依存性}

スーパークラスは子要素の出現順に依存しない．
\begin{equation}
(A, B) = (B, A)
\end{equation}

\paragraph{性質2：null の透過性}

スーパークラスの子要素としての $\mathrm{null}$ は，存在しないことと等価である．
\begin{equation}
(A, \mathrm{null}, B) = (A, B)
\end{equation}

\paragraph{性質3：逆演算の存在}

スーパークラス $\mathcal{I}$ に対する任意の演算について，その逆演算が存在する．

上記の性質1〜3を満たす演算を $*$ と定める．

\paragraph{メタ情報の定義}

スーパークラスの基底クラスは，情報 $I$ そのものとは区別される\textbf{メタ情報}を持つ．メタ情報は以下の2つから構成される：

\begin{itemize}
  \item \textbf{カウンタ $n$}：ディレイ回路の通過回数を積算する．初期値は $\mathrm{null}$ とする．
  \item \textbf{演算ファンクション $*$}：ディレイ回路を通過する際に実行される演算．初期値は $\mathrm{null}$ とする．演算ファンクション $*$ は $n$ を参照することができる．
\end{itemize}

メタ情報は情報 $I$ に付随するが，$I$ の値そのものではない．したがって，前稿~\cite{ref1} \S 2.2 における情報 $I$ の定義は変更されない．

$n$ は逆演算の例外であり，逆演算を行っても保存される．

\subsection{$F_D$ の拡張定義}

前稿~\cite{ref1} \S 2.2 で定義された $F_D$ の演算規則（入力情報を内部情報にコピーし，出力に渡す）は変更しない．スーパークラスを通過する場合，$F_D$ は上記の演算規則に加えて，通過する情報のメタ情報に対して以下の操作を行う：

\begin{enumerate}
  \item カウンタ $n$ を $+1$ する（$n = \mathrm{null}$ の場合は $0$ から $1$ にする）
  \item 演算ファンクション $*$ が $\mathrm{null}$ でなければ実行する
\end{enumerate}

メタ情報の操作主体は $F_D$ であり，情報 $I$ 自身が「$F_D$ を通過した」と認識する必要はない．

NOT演算回路はメタ情報の操作を行わない．これにより，前稿~\cite{ref1} \S 2.3 で「NOT演算回路の演算は遅延を伴わないため，ディレイ回数 $n$ には含めない」と述べた事実が，構造的に裏付けられる．

自身で子クラスを生成する演算の場合，$n$ は初期化できる．

\subsection{スーパークラスとディレイ回路の整合性}

スーパークラス $\mathcal{I}$ は，前稿~\cite{ref1} \S 2.2 における情報 $I$ の定義（「数値表現できるあらゆる情報」）を満たす．すなわち，スーパークラスに属する情報の組 $(A, B)$ や $(C_1, C_2, \ldots, C_m)$ は，それぞれ一つの情報 $I$ として扱うことができる．

したがって，ディレイ回路 $F_D$ および NOT演算回路の定義（\cite{ref1} \S 2.2〜\S 2.4）は，変更なくスーパークラスの情報に適用される．\S 3.3 で追加されたメタ情報の操作は，$F_D$ の元の演算規則の変更ではなく，新たな操作層の追加である．

\subsection{系の外部からの介入}

前稿~\cite{ref1}で定義された $F_D$ の演算規則および \S 3.3 のメタ情報の操作は，系の内部における決定論的な動作を規定する．本節では，系の外部からの介入を定義する．

\textbf{外部介入}：系の外部から，特定の $F_D$ の内部情報を強制的に上書きする操作を許す．この操作は系の内部の演算規則とは独立に行われる．

外部介入は \S 3.2 性質3 の例外であり，外部介入を受けた情報は逆演算の対象外となる．すなわち，外部介入は系の内部に不可逆性を持ち込む．

系の内部観測者から見れば，外部介入は因果律の外から突然現れる事象である．

注：本定義は，本モデルが階層構造を許容することを示唆する．上位の系と下位の系の関係については本論の射程外とする．

\section{考察}

\section{結論}

本モデルの方程式化や物理学的フレームワークへのマッピングは本論の射程外である．

\begin{thebibliography}{9}
\bibitem{ref1} 木原範昭「情報伝達の情報論的整理」研究ノート，2026年4月．
\end{thebibliography}

\end{document}
